\section{Quadratic identities}\label{sec:04-quadratic}

The entropy estimate reduces the cone inequality to $\mathcal E\le0$.
We now rewrite the signed term $\delta[d]-\tau_T H_v$ in
$\mathcal E$. We first express the covariance of the layer integrals
$H_f$ through a scalar kernel. Comparing that kernel with its Gaussian
reference gives identities for both $\delta[l]$ and $\delta[l^2]$.
We then apply them to $d=F-|u|^2$, separating the functionals of the
linear field $L$ from the residual field $R=A-L$ and the layer defects.

Throughout this section the choices in \eqref{eq:8} are fixed.
In particular, \(T,\Sigma=I+T\), and the coefficients \(M_i\) of
\(L=\sum_iG_iM_i\) are deterministic. Scalar tests are smooth, with
their derivatives of at most polynomial growth, as in the preceding
sections.

\subsection{Gaussian covariance with a deterministic matrix weight}

Let \(\mathcal N_G\) be the nonnegative Gaussian number operator in
\(G\in\mathbb R^m\), given on smooth fields by
\(\mathcal N_G=G\cdot\nabla_G-\Delta_G\) and acting entrywise on
matrix fields, and put
\[
 \mathcal K_G=(1+\mathcal N_G)^{-1}.
\]
For symmetric matrix fields \(E,F\), recall that
\[
 \langle E,F\rangle_\Sigma
   =\frac1m\mathbb E\operatorname{tr}\bigl(\Sigma(EF+FE)/2\bigr)
   =\frac1m\mathbb E\operatorname{tr}(\Sigma EF).
\]
The second equality follows by transposition and cyclicity of trace;
it does not require commutation. Since \(\Sigma>0\), this is a
positive inner product. Its Gaussian form norm is
\[
 \|E\|_{1,\Sigma}^2
   :=\langle E,(1+\mathcal N_G)E\rangle_\Sigma
    =\tau_\Sigma E^2+
       \sum_{i=1}^m\tau_\Sigma(\partial_{G_i}E)^2,
\]
with the left side interpreted as a closed quadratic form.

The covariance formula below is the finite-dimensional Gaussian form of
the Ornstein--Uhlenbeck semigroup representation in
\cite[Proposition~2.1, Equation~(2.4)]{HoudrePrivault2002}.
We include its Hermite proof in the normalization used here.

\begin{lemma}[Gaussian covariance formula]\label{lem:gaussian-covariance}
 Let \(\mathcal X,\mathcal Y\) be centered elements of the Gaussian
 Sobolev space \(W^{1,2}(\mathbb R^m;\mathbb R^m)\). Then
 \[
  \mathbb E(\mathcal X\cdot\mathcal Y)
    =\mathbb E\sum_{i=1}^m
       (\partial_{G_i}\mathcal X)\cdot
       \mathcal K_G(\partial_{G_i}\mathcal Y).
 \]
 Moreover, for every square-integrable symmetric matrix field \(E\),
 \(\mathcal K_GE\) belongs to the form domain of
 \(1+\mathcal N_G\), for the deterministic weight \(\Sigma\).
\end{lemma}

\begin{proof}
Use the orthonormal multivariate Hermite basis, indexed by
\(\mathbf a=(a_1,\ldots,a_m)\in\{0,1,2,\ldots\}^m\); see, for example,
\cite[Sections 5--6]{Bell2015}. The number operator has eigenvalue
\(|\mathbf a|=\sum_i a_i\), and differentiation in coordinate \(i\)
lowers the index by \(e_i\) with factor \(\sqrt{a_i}\). Thus, on a
nonconstant mode, differentiation and application of \(\mathcal K_G\)
give the multiplier
\[
 \sum_{i:a_i>0}\frac{a_i}{1+|\mathbf a|-1}=1.
\]
This proves the covariance formula for finite Hermite sums and then
for \(W^{1,2}\) fields by convergence of the coefficients and their
weak derivatives.

For a matrix field with coefficient \(E_{\mathbf a}\), the
contribution of \(\mathcal K_GE\) to its squared form norm is
\[
 \frac{\langle E_{\mathbf a},E_{\mathbf a}\rangle_\Sigma}
      {1+|\mathbf a|}.
\]
It is summable whenever \(E\) is square-integrable. The deterministic
weight \(\Sigma\) acts only on matrix coefficients, so it commutes
with the Hermite degree decomposition. The same expansion gives
the stated form norm and its Dirichlet expression.
\end{proof}

The regularity needed here is that of the vector fields in
Lemma~\ref{lem:projection-sobolev}. We will not differentiate the
measurable projection derivatives \(P_z\), or the fields \(H_f\).
The latter need only be square-integrable before applying
\(\mathcal K_G\).

Define the scalar bilinear kernel
\[
 \mathcal Q[f,j]
  =\iint f(x)j(z)\,p(\min(x,z))r_1(\max(x,z))\,dx\,dz.
\]
Its density is nonnegative and has both marginals \(\mu\):
integrating at a fixed coordinate \(z\) gives
\(a(z)r_1(z)+p(z)B(z)=h'(z)\), by \eqref{eq:6}.

The covariance also records the contact between projections at different
layers. Define
\[
 \operatorname{ct}(x,z)
 =\frac1m\mathbb E[X_{\min(x,z)}\cdot Y_{\max(x,z)}],
 \qquad
 \operatorname{Ct}(f,j)
 =\iint\operatorname{ct}(x,z)f'(x)j'(z)\,dx\,dz.
\]
Positive duality gives $\operatorname{ct}\ge0$. At a single layer,
Moreau's orthogonality gives $X_z\cdot Y_z=0$; at distinct layers
the contact need not vanish. These integrals converge absolutely:
on $x\le z$, the same three-region argument used for
$\operatorname{nt}$ applies to $X_x$ and $Y_z$, with the remaining
factors controlled by their polynomial moment bounds.

\begin{lemma}[Covariance of the layer integrals]\label{lem:layer-covariance}
 For smooth scalar tests \(f,j\) with \(\gamma f=\gamma j=0\),
 \begin{equation}\tag{13}\label{eq:13}
  \langle H_f,\mathcal K_G H_j\rangle_\Sigma
    =\mathcal Q[f,j]+\operatorname{Ct}(f,j).
 \end{equation}
\end{lemma}

\begin{proof}
Consider the centered vector fields
\[
 \mathcal W_f
    =\int_{\mathbb R}
       \bigl(p(z)Z-X_z+\mathbb E X_z\bigr)f'(z)\,dz
\]
and \(\mathcal W_j\). By Lemma~\ref{lem:projection-sobolev}, they
belong to Gaussian \(W^{1,2}\), with Jacobians
\(H_f\Sigma^{1/2}\) and \(H_j\Sigma^{1/2}\). Applying
Lemma~\ref{lem:gaussian-covariance} and dividing by \(m\) gives
\(\langle H_f,\mathcal K_GH_j\rangle_\Sigma\).

We also compute this covariance from the layers. If \(x\leq z\),
the identities \(X_z=Z+\xi_z+Y_z\) and
\(\xi_z=a(z)U-\mathbb E Y_z\) imply
\[
 \frac1m\mathbb E
   \bigl[(X_x-\mathbb E X_x)\cdot(X_z-\mathbb E X_z)\bigr]
   =h(x)-a(x)B(z)+\operatorname{ct}(x,z).
\]
Indeed, the terms involving \(a(x)a(z)|U|^2\) cancel after subtracting
the product of the means. For the centered differences defining
\(\mathcal W_f\), the remaining covariance kernel, apart from
\(\operatorname{ct}(x,z)-a(\min(x,z))B(\max(x,z))\), is
\[
 s_0p(x)p(z)-p(x)h(z)-h(x)p(z)+h(\min(x,z)).
\]
Since \(\mu\) has mass \(s_0\) and cumulative function \(h\), this
equals
\[
 \int
  \bigl(p(x)-\mathbf1_{\{y\leq x\}}\bigr)
  \bigl(p(z)-\mathbf1_{\{y\leq z\}}\bigr)\,d\mu(y).
\]
The centered-test identity
\[
 \int\bigl(p(z)-\mathbf1_{\{y\leq z\}}\bigr)f'(z)\,dz=f(y)
\]
therefore turns its integral against \(f'(x)j'(z)\) into \(\mu[fj]\).

The mixed distributional derivative of
\(a(\min(x,z))B(\max(x,z))\) is
\[
 d\mu(x)\,\delta_x(dz)
  -p(\min(x,z))r_1(\max(x,z))\,dx\,dz.
\]
Off the diagonal this follows from \(a'=p\) and \(B'=-r_1\).
On the diagonal, the jump is \(pB+a r_1=h'\).
Integrating by parts in both variables thus gives
\(\mu[fj]-\mathcal Q[f,j]\) for this kernel. Subtraction proves
\eqref{eq:13}.

For noncompact tests, insert smooth cutoffs and pass to the limit.
The tail estimates imply integrability with polynomial weights of
\(a(\min(x,z))B(\max(x,z))\): one separates the ordered region
into \(x\leq z\leq0\), \(x\leq0\leq z\), and
\(0\leq x\leq z\). The diagonal measure and the density of
\(\mathcal Q\) have all polynomial moments, since their marginals
are \(\mu\). These bounds justify the integrations by parts and the
limit.
\end{proof}

\subsection{The scalar reference kernel and its variation}

Let \(\mathcal Q_0\) be the same kernel with \(r_1=q_0=1-p\).
The following identity converts its variation into the functional
\(\delta\) from \eqref{eq:9}.

\begin{lemma}[Reference inversion and variation]\label{lem:kernel-variation}
 Let \(h_f^*,h_j^*\) be smooth scalar tests, and set
 \(f=(1+\mathcal N)h_f^*\), \(j=(1+\mathcal N)h_j^*\).
 Then
 \begin{equation}\tag{14}\label{eq:14}
 \begin{aligned}
  \mathcal Q_0[f,j]&=\gamma[f h_j^*],\\
  (\mathcal Q-s_0\mathcal Q_0)[f,j]
    &=\Delta_\mu w+\delta[(h_f^*)'(h_j^*)'],&
  w'&=f'h_j^*+j'h_f^* .
 \end{aligned}
 \end{equation}
 The choice of the additive constant in \(w\) is immaterial.
\end{lemma}

\begin{proof}
Write \(b_1=h_f^*\), \(b_2=h_j^*\). Direct differentiation gives
\[
 \int_{-\infty}^z p f=a b_1-p b_1',\qquad
 \int_z^\infty q_0 f=B_0 b_1+q_0 b_1',
\]
and the analogous formulas for \(j\). Boundary terms vanish by
polynomial growth and the Gaussian tails. As
\(q_0a+pB_0=\phi\), insertion into \(\mathcal Q_0\) gives
\(\mathcal Q_0[f,j]=\gamma[jb_1]=\gamma[fb_2]\), using
self-adjointness of \(1+\mathcal N\) under \(\gamma\).

Put
\[
 H_1=fb_2+jb_1,\qquad J_1=fb_2'+jb_1',\qquad w'=H_1'-J_1.
\]
Split the kernel variation into its two ordered half-planes. Symmetry
and the preceding primitive formulas give
\[
 \begin{aligned}
 (\mathcal Q-s_0\mathcal Q_0)[f,j]
 &=\int\Delta r(z)\left[
       j(z)\int_{-\infty}^z p(x)f(x)\,dx
       +f(z)\int_{-\infty}^z p(x)j(x)\,dx\right]dz\\
 &=\int\Delta r(aH_1-pJ_1).
 \end{aligned}
\]
Use \(\Delta r=-(\Delta B)'\),
\(\Delta_\mu w=-\int h_\delta w'\), and
\((a\Delta B+h_\delta)'=2p\Delta B\). After integration by parts,
subtracting \(\Delta_\mu w\) leaves
\[
 -\int\bigl[p\Delta B\,H_1+(p\Delta r+h_\delta)J_1\bigr].
\]
The product rule for \(\mathcal N\) gives
\[
 \begin{aligned}
 H_1&=(2+\mathcal N)(b_1b_2)+2b_1'b_2',\\
 J_1&=(b_1b_2)'+2z\,b_1'b_2'-(b_1'b_2')'.
 \end{aligned}
\]
Equation~\eqref{eq:12} cancels the terms involving \(b_1b_2\).
The remaining expression is
\(\delta[b_1'b_2']\) by \eqref{eq:9}, proving \eqref{eq:14}.
All boundary terms contain the weighted differences controlled by
the preceding tail estimates. Finally \(\Delta_\mu\) has total
mass zero, so constants in \(w\) do not affect the formula.
\end{proof}

\subsection{Defects from spectral thresholds}

For a deterministic symmetric weight $\Lambda$, extend the kernel
$\operatorname{nt}$ of Section~\ref{sec:03-entropy} by
\[
 \begin{aligned}
 \operatorname{nt}_\Lambda(x,z)
 &=\tau_\Lambda\bigl(P_{\min(x,z)}\circ(I-P_{\max(x,z)})\bigr),\\
 \operatorname{Nt}_\Lambda(f,j)
 &=\iint\operatorname{nt}_\Lambda(x,z)f'(x)j'(z)\,dx\,dz,
 \qquad N_\Lambda=\operatorname{Nt}_\Lambda(\iota,\iota).
 \end{aligned}
\]
We omit the subscript $I$. The layer-tail argument again gives
absolute convergence for smooth tests with polynomially bounded
derivatives. A general weighted kernel need not be nonnegative.

For a symmetric random matrix \(B_*\) with all finite moments, define
\[
 \Theta_z(B_*)=\mathbf1_{(-\infty,z]}(B_*),\qquad
 E_z^{B_*}=P_z-\Theta_z(B_*).
\]
We will use \(B_*=A\) and \(B_*=L\). For a deterministic symmetric
matrix \(\Lambda\), put
\[
 e_\Lambda^{B_*}(f,j)
  =\int f'(z)
      \langle E_z^{B_*},j(B_*)-j(z)I\rangle_\Lambda\,dz.
\]
An omitted superscript means \(B_*=L\). An omitted trace weight means
\(\Lambda=I\).

Define a measure on layer and spectral-endpoint pairs by
\[
 \int F_0(z,x)\,d\beta^{B_*}(z,x)
   =\int dz\,\tau\!\left[
       E_z^{B_*}(B_*-zI)F_0(z,B_*)\right].
\]
This is a nonnegative measure. In a spectral basis of \(B_*\), with
eigenvalues \(x_i\), its density is the normalized expected sum of
\((E_z^{B_*})_{ii}(x_i-z)\). If \(x_i>z\), the first factor is
\((P_z)_{ii}\geq0\); if \(x_i\leq z\), it is
\((P_z)_{ii}-1\leq0\). Its mass is
\[
 \mathcal B^{B_*}=e^{B_*}(\iota,\iota).
\]
Write \(\beta,\mathcal B\) for \(B_*=L\).
The layer tails, H\"older's inequality, and the moments of \(B_*\)
give convergence of these integrals with polynomial weights.
For instance, the spectral thresholds have arbitrarily high inverse
polynomial tail bounds from the moments of \(B_*\), while the
projection-layer errors have the stronger Gaussian tail bounds
already established.

Set
\[
 M_f^{B_*}
    =H_f-f(B_*)+(\gamma f)I
    =-\int E_z^{B_*}f'(z)\,dz.
\]
These matrix fields are distinct from the deterministic linear
coefficients \(M_i\) in \eqref{eq:7}.

\begin{lemma}[Quadratic threshold defect]\label{lem:threshold-defect}
 For every such \(B_*,f,\Lambda\),
 \begin{equation}\tag{15}\label{eq:15}
  \langle M_f^{B_*},M_f^{B_*}\rangle_\Lambda
    =2e_\Lambda^{B_*}(f,f)-\operatorname{Nt}_\Lambda(f,f).
 \end{equation}
\end{lemma}

\begin{proof}
Abbreviate \(\Theta_y=\Theta_y(B_*)\). For \(x<z\),
\(\Theta_x\Theta_z=\Theta_x\), and expansion gives
\[
 E_x^{B_*}\circ E_z^{B_*}+P_x\circ(I-P_z)
   =P_x\circ(I-\Theta_z)+(I-P_z)\circ\Theta_x.
\]
Multiply by \(f'(x)f'(z)\), integrate over \(x<z\), and include the
opposite order. The first two terms give the left side of
\eqref{eq:15} and \(\operatorname{Nt}_\Lambda(f,f)\).
The terms on the right give \(2e_\Lambda^{B_*}(f,f)\), by the
spectral identity
\[
 f(B_*)-f(y)I
   =\int\bigl(\mathbf1_{\{y\leq l\}}I-\Theta_l\bigr)f'(l)\,dl.
\]
All rearrangements are justified by the polynomially weighted
integrability above. In particular, no nesting of \(P_z\), or
commutation of \(P_z\) with a threshold or with \(\Lambda\), has
been assumed.
\end{proof}

For the linear test \(\iota(x)=x\), we have
\(M_\iota^L=A-L=R\) and \(M_\iota^A=0\). Thus
\[
 2e_\Lambda(\iota,\iota)=\tau_\Lambda R^2+N_\Lambda,\qquad
 2e_\Lambda^A(\iota,\iota)=N_\Lambda.
\]
The next identities compare a layer integral with evaluation at \(L\).
Use
\[
 \Delta_\Lambda j=\tau_\Lambda\bigl(j(L)-(\gamma j)I\bigr).
\]

\begin{lemma}[Layer-to-spectral comparison]\label{lem:layer-spectral}
 For smooth scalar tests \(f,j\), with \(\gamma f=0\) in the last
 identity,
 \begin{equation}\tag{16}\label{eq:16}
 \begin{aligned}
  \tau_\Lambda H_j
   &=\Delta_\Lambda j+\langle R,j'(L)\rangle_\Lambda
                         +e_\Lambda(\iota,j'),\\
  \tau_\Sigma H_j&=\Delta_\mu j,\\
  \langle H_f,j(L)\rangle_\Sigma
   &=\mu[fj]+(\Delta_\Sigma-\Delta_\mu)w-e_\Sigma(f,j),
       \qquad w'=fj'.
 \end{aligned}
 \end{equation}
\end{lemma}

\begin{proof}
The first identity follows from \(R=-\int E_z^L\,dz\):
\[
 \langle R,j'(L)\rangle_\Lambda+e_\Lambda(\iota,j')
     =-\int j'(z)\tau_\Lambda E_z^L\,dz.
\]
For the second, integration by parts in
\(\tau_\Sigma H_j=\int(s_0p-h)j'\) gives
\(\Delta_\mu j\).
For the third, subtract \(f(L)\) from \(H_f\) in the pairing and
write \(j(L)=(j(L)-j(z)I)+j(z)I\).
The first part gives \(-e_\Sigma(f,j)\). For the scalar part use
\(f'j=(fj-w)'\) and the second identity. Combining with
\(\tau_\Sigma(fj)(L)\) gives the last line of \eqref{eq:16}.
\end{proof}

\subsection{Linear and quadratic profile identities}

In the reference-variation formula~\eqref{eq:14}, the derivatives of
the two inverse tests appear as the product inside $\delta$.
Using the same centered primitive of $l$ in both slots gives
$\delta[l^2]$; pairing it with
$(1+\mathcal N)^{-1}\iota=\iota/2$ gives $\delta[l]/2$.
These are the two choices used below.
For a smooth scalar test $l$, let $h_l$ be its antiderivative
centered under the one-dimensional Gaussian. Apply $1+\mathcal N$
to this primitive to define $f_l$; its derivative will be denoted
by $S_l$. Thus
\[
 h_l'=l,\qquad \gamma h_l=0,\qquad
 f_l=(1+\mathcal N)h_l,\qquad S_l=f_l'=(2+\mathcal N)l,
 \qquad \psi_l'=(\mathcal N-1)h_l.
\]
The last definition introduces the additional primitive \(\psi_l\)
that will occur in the scalar variation term. Choose its additive
constant linearly in \(l\). Gaussian integration by parts shows that
\(f_l\) is centered under \(\gamma\), so the layer covariance formula
applies to it.

We now compare the two matrix fields obtained from \(h_l\): the
resolvent \(\mathcal K_GH_{f_l}\) of its layer test, and the direct
evaluation \(h_l(L)\). Denote their difference by \(d_l^*\), its
squared form norm by \(D_\Sigma(l)^2\), and its pairing with the
remainder \(R\) by \(\sigma_\Sigma(l)\). We also define the difference
\(J_\Sigma(l)\) between the displayed spectral term and the Dirichlet
term of \(h_l(L)\):
\[
 \begin{aligned}
 d_l^*&=\mathcal K_GH_{f_l}-h_l(L),&
 D_\Sigma(l)^2&=\|d_l^*\|_{1,\Sigma}^2,&
 \sigma_\Sigma(l)&=2\langle R,d_l^*\rangle_\Sigma,\\
 J_\Sigma(l)&=\tau_\Sigma\bigl(Lh_{l^2}(L)\bigr)
       -\sum_{i=1}^m\tau_\Sigma
                   \bigl(\partial_{G_i}h_l(L)\bigr)^2 .
 \end{aligned}
\]
For vector-valued \(l\), sum \(D_\Sigma(l)^2\) and \(J_\Sigma(l)\)
over its components.
The quadratic identity below separates the nonnegative quantity
\(D_\Sigma(l)^2\) from \(J_\Sigma(l)\), which depends only on \(L\)
and the deterministic trace weight.

These form quantities are finite. In a spectral basis of \(L\),
\[
 \bigl(\partial_{G_i}h_l(L)\bigr)_{ab}
   =h_l^{[1]}(x_a,x_b)(M_i)_{ab},\qquad
 h_l^{[1]}(x,y)=
 \begin{cases}
  \dfrac{h_l(x)-h_l(y)}{x-y},&x\ne y,\\
  h_l'(x),&x=y.
 \end{cases}
\]
The formula follows for polynomials by differentiating products.
Approximation of a smooth function and its first derivative on compact
intervals gives the general formula; the Hilbert--Schmidt norm of
this derivative is bounded by the supremum of the scalar derivative
on the spectral interval times the norm of its matrix direction.
Polynomial growth and the Gaussian moments of the linear matrix
\(L\) then give square-integrable derivatives.
Lemma~\ref{lem:gaussian-covariance} supplies the form regularity of
\(\mathcal K_GH_{f_l}\).

\begin{lemma}[Profile identities]\label{lem:profile-quadratics}
 For every smooth scalar test \(l\),
 \begin{equation}\tag{17}\label{eq:17}
 \begin{aligned}
  \delta[l]
   &=\Delta_\Sigma\psi_l+\sigma_\Sigma(l)
      -2e_\Sigma(\iota,h_l)-e_\Sigma(f_l,\iota)
      -2\operatorname{Ct}(\iota,f_l),\\
  D_\Sigma(l)^2+J_\Sigma(l)
   &=\delta[l^2]-\Delta_\Sigma\psi_{l^2}
      +2e_\Sigma(f_l,h_l)+\operatorname{Ct}(f_l,f_l).
 \end{aligned}
 \end{equation}
\end{lemma}

\begin{proof}
For the linear identity, evaluate
\(\mu[\iota f_l]-2\mathcal Q[\iota,f_l]\) in two ways.
Since \((1+\mathcal N)^{-1}\iota=\iota/2\), Equation~\eqref{eq:14}
gives
\[
 \mu[\iota f_l]-2\mathcal Q[\iota,f_l]
    =\Delta_\mu\psi_l-\delta[l].
\]
The reference Gaussian term vanishes by self-adjointness of
\(1+\mathcal N\), and the derivative of the remaining primitive is
\(f_l-2h_l=\psi_l'\).

For the second evaluation, replace \(\mathcal Q[\iota,f_l]\) using
Equation~\eqref{eq:13} and replace \(\mu[\iota f_l]\) using the
last line of Equation~\eqref{eq:16}. Since \(H_\iota=A=L+R\)
and \(\mathcal K_GL=L/2\), the same quantity
\(\mu[\iota f_l]-2\mathcal Q[\iota,f_l]\) equals
\[
 -(\Delta_\Sigma-\Delta_\mu)v_0+e_\Sigma(f_l,\iota)
   -2\langle R,\mathcal K_GH_{f_l}\rangle_\Sigma
   +2\operatorname{Ct}(\iota,f_l),\qquad v_0'=f_l.
\]
Equate these two evaluations. Equation~\eqref{eq:16}, applied to
\(\psi_l-v_0\), whose derivative is \(-2h_l\), rewrites the remaining
difference between \(\Delta_\mu\) and \(\Delta_\Sigma\).
Its pairing with \(R\) combines with
\(2\langle R,\mathcal K_GH_{f_l}\rangle_\Sigma\) to give
\(2\langle R,d_l^*\rangle_\Sigma=\sigma_\Sigma(l)\).
The other terms are precisely the first line of Equation~\eqref{eq:17}.

For the quadratic identity write \(b_1=h_l\), \(f=f_l\).
Expand the squared form norm of
\(d_l^*=\mathcal K_GH_f-b_1(L)\).
In Hermite degree \(j\), the form contributes a factor \(1+j\)
and \(\mathcal K_G\) contributes its reciprocal. Thus the cross term
is \(-2\langle H_f,b_1(L)\rangle_\Sigma\), and the square of
\(\mathcal K_GH_f\) is
\(\langle H_f,\mathcal K_GH_f\rangle_\Sigma\). This uses the form
regularity already proved, not a derivative of \(H_f\). Consequently
\[
 \begin{split}
 D_\Sigma(l)^2
  ={}&\langle H_f,\mathcal K_GH_f\rangle_\Sigma
       -2\langle H_f,b_1(L)\rangle_\Sigma
       +\tau_\Sigma b_1(L)^2\\
     &+\sum_i\tau_\Sigma
                    \bigl(\partial_{G_i}b_1(L)\bigr)^2 .
 \end{split}
\]
Adding \(J_\Sigma(l)\) cancels exactly the same weighted Dirichlet
term. By \eqref{eq:13}, the remaining expression is
\[
 \mathcal Q[f,f]+\operatorname{Ct}(f,f)
   -2\langle H_f,b_1(L)\rangle_\Sigma
   +\tau_\Sigma\bigl(b_1(L)^2+Lh_{l^2}(L)\bigr).
\]
Insert \eqref{eq:14} for \(\mathcal Q[f,f]\) and
\eqref{eq:16} for the pairing. The reference terms cancel because
\[
 \gamma[xh_{l^2}]=\gamma[l^2]
    =\gamma[(f-b_1)b_1].
\]
The \(\Delta_\mu\) terms have the primitive
\[
 2\int^x f'b_1-2fb_1+2\int^x fb_1',
\]
whose derivative is zero. Besides
\(\delta[l^2]+\operatorname{Ct}(f,f)+2e_\Sigma(f,b_1)\),
what remains is \(\Delta_\Sigma\) applied to
\[
 b_1^2+xh_{l^2}-2\int^x fb_1'.
\]
Its derivative is
\[
 h_{l^2}-x l^2+2ll'
    =-(\mathcal N-1)h_{l^2}=-\psi_{l^2}'.
\]
Since \(\Delta_\Sigma\) annihilates constants, this proves the second
identity. All form expansions used a deterministic weight; none
requires it to commute with the matrix fields.
\end{proof}

\subsection{Combining the fixed profiles}

We can now arrange the signed term from the entropy estimate.
For a two-variable layer profile \(\Phi(z,x)\), define
\[
 \mathcal L_\Lambda[\Phi]
  =\int dz\,
    \left\langle E_z^L,\,
       \left.x\longmapsto\int_z^x\Phi(z,y)\,dy\right|_{x=L}
    \right\rangle_\Lambda .
\]
With
\(\overline\Phi(z,x)=\int_0^1\Phi(z,z+t(x-z))\,dt\), this means
\[
 e_\Lambda(f,j)=\mathcal L_\Lambda[f'(z)j'(x)],
 \qquad
 \mathcal L_I[\Phi]=\int\overline\Phi(z,x)\,d\beta(z,x).
\]
This convention includes coincident endpoints and either ordering
of the endpoints.

\begin{proposition}[Decomposition of the entropy term]\label{prop:quadratic-decomposition}
 With the fixed profiles \(F=d+|u|^2=c+2kq\), one has
 \begin{equation}\tag{18}\label{eq:18}
 \begin{aligned}
  \delta[d]-\tau_T H_v
  \leq{}&\Delta_{\rm p}+\sigma_\Sigma(F)
       -D_\Sigma(u)^2-J_\Sigma(u)-\langle R,v'(L)\rangle_T\\
       &-\mathcal L_I[H]-\mathcal L_T[H+v''(x)],\\
  \Delta_{\rm p}
    ={}&\Delta_I\psi_d+\Delta_T(\psi_d-v).
 \end{aligned}
 \end{equation}
\end{proposition}

\begin{proof}
Use the first line of \eqref{eq:17} for \(F\), subtract the second
line for each component of \(u\), and use
\(\delta[d]=\delta[F]-\delta[|u|^2]\).
The combined layer profile for the weight \(\Sigma\) is
\[
 2F(x)+S_F(z)-2S_u(z)\cdot u(x)
 =4c+4kq(x)+2(k-q(x))S_q(z)-2S_\pi(z)\cdot\pi(x)
 =H(z,x).
\]
The contact integral being subtracted is
\[
 \iint \operatorname{ct}(x,z)
   \bigl(S_F(x)+S_F(z)-S_u(x)\cdot S_u(z)\bigr)\,dx\,dz.
\]
It is nonnegative: \(\operatorname{ct}\geq0\), and its scalar
factor is the second positive expression in \eqref{eq:3}.
Discarding this integral preserves the upper bound.

Finally, subtract \(\tau_TH_v\) using the first line of
\eqref{eq:16}. Its derivative correction is
\(-\langle R,v'(L)\rangle_T\), and its layer correction is
\(-e_T(\iota,v')=-\mathcal L_T[v''(x)]\).
Splitting \(\Sigma=I+T\) gives the two layer terms in
\eqref{eq:18}, while
\(\Delta_\Sigma\psi_d-\Delta_Tv=\Delta_{\rm p}\).
\end{proof}

Equation~\eqref{eq:18} leaves two kinds of error: the field \(R\) and
the defects \(E_z^L\). The next section bounds them while retaining a
nonnegative layer remainder. The terms \(\Delta_{\rm p}\) and
\(J_\Sigma(u)\), which depend only on the linear field \(L\), will
then be evaluated spectrally.
